\documentclass[12pt,a4paper]{article}
\usepackage[margin=1in]{geometry}
\usepackage{enumitem}
\usepackage{booktabs}
\usepackage{longtable}
\usepackage{hyperref}
\usepackage{xcolor}

\hypersetup{colorlinks=true, linkcolor=blue!60!black, urlcolor=blue!60!black}

\title{Paper Status Assessment and Recommendations\\[6pt]
\large ``Quantity-Space Prior Specification for Bayesian Hierarchical Models''}
\author{Technical Note}
\date{February 2026}

\begin{document}
\maketitle

\section{Overview}

This note provides a section-by-section assessment of the current draft,
identifies issues ranging from minor to structural, and recommends an order of
operations for bringing the paper to submission quality. The assessment is based
on a full read of \texttt{doc/main.tex} as of February 25, 2026.

The paper is in strong shape methodologically. Sections~1--5 contain the full
argument for the primary application (RDU with real data), and the methodology
sections (3--4) are tight after recent editing. The main gaps are the two stub
sections (CPT and Discounting simulation studies), and a collection of
presentation issues distributed across the existing sections.


\section{Section-by-Section Assessment}

\subsection{Abstract (lines 86--112)}

The abstract is well-written and accurately summarizes both contributions. Two
observations:

\begin{enumerate}[label=(\alph*)]
  \item The phrase ``reduces the remaining degrees of freedom to two
    interpretable parameters'' is slightly ambiguous---it could be read as
    reducing from 4$p$ to 2 total, when the actual reduction is from 4$p$ to
    2$p$ + 2. The conclusion states this correctly; the abstract should match.
  \item The abstract mentions ``any distributional family'' for the hyperprior.
    This is true for the moment-matching framework but the worked examples are
    all NIG. Consider softening to ``demonstrated for standard conjugate
    families and extensible to others'' or leaving as-is and ensuring the
    Skew-Normal appendix is cited as evidence.
\end{enumerate}


\subsection{Section 1: Introduction (lines 128--213)}

Solid. The two-camp framing is clear. One remaining issue:

\begin{enumerate}[label=(\alph*)]
  \item[{[x]}] Line 167: ``an quantity-space'' --- fixed to ``a quantity-space.''
  \item The road map paragraph (lines 201--213) references Sections 6 and 7 but
    does not warn the reader that these are forthcoming simulation studies. If
    the paper will be submitted with stubs, the road map should acknowledge
    this. If the stubs will be filled in before submission, no change is needed.
\end{enumerate}


\subsection{Section 2: Literature --- Two Camps (lines 216--378)}

This section is long (approximately 4 pages of double-spaced text) and
currently the weakest section in terms of focus. The belief-elicitation
literature (lines 303--338) is interesting but tangential: the paper does not
actually elicit beliefs from subjects. This subsection reads as a detour that a
reviewer might ask to be cut.

\begin{enumerate}[label=(\alph*)]
  \item Consider condensing the belief-elicitation material to 2--3 sentences
    acknowledging the literature, or moving it to a footnote. The key point---
    that translating domain knowledge into priors is hard---is already made by
    the Kadane and Gabry citations.
  \item The transition paragraph (lines 363--378) does important work connecting
    Segert, Bland, and Manderson \& Goudie to the present contribution. This
    is the strongest part of the section and should be preserved.
  \item Camp~1 (lines 239--288) is well-balanced. No changes recommended.
\end{enumerate}


\subsection{Section 3: General Hyperprior Specification (lines 382--611)}

This section has been recently edited and is in good shape. Remaining issues:

\begin{enumerate}[label=(\alph*)]
  \item The section title says ``General Hyperprior Specification from Generated
    Quantities'' but the section covers \emph{two} contributions: (i) the
    Jacobian correction for quantity-space priors (\S3.1) and (ii) the variance
    decomposition (\S3.2--3.3). Consider whether the title should reflect both,
    or whether Section~3 should be split into two sections. If the paper is
    targeting 9 sections, splitting Section~3 into ``Quantity-Space Prior
    Specification'' (current \S3.1) and ``Hyperprior Construction via Variance
    Decomposition'' (current \S3.2--3.3) would create a cleaner structure and
    better match the ``two contributions'' framing. The section numbers in the
    road map would then shift by one.
  \item Line 421: ``we will only focus'' --- the paper otherwise uses first
    person singular (``I''). This should be ``I focus'' or ``the paper focuses.''
  \item The ``no free lunch'' paragraph (lines 597--611) is important but its
    placement at the end of \S3.3 makes it feel like an afterthought. It might
    work better as the opening of a brief ``Summary'' paragraph or as the
    closing of \S3.2, immediately after the law of total variance is introduced.
\end{enumerate}


\subsection{Section 4: Application --- RDU with Real Data (lines 614--1078)}

This is the longest section and carries the empirical weight of the paper. It
is well-structured but has several issues.

\subsubsection{Model specification (\S4.1, lines 620--661)}

\begin{enumerate}[label=(\alph*)]
  \item The model specification (RDU, CRRA, Prelec, Fechner) is standard and
    clear.
  \item The Wilcox (2011) contextual utility normalization is mentioned but not
    explained in detail. A reviewer unfamiliar with this could stumble. A
    one-sentence clarification or a forward reference to the appendix would
    help.
\end{enumerate}

\subsubsection{CE mapping and injectivity (\S4.2--4.3, lines 662--798)}

\begin{enumerate}[label=(\alph*)]
  \item The CE-space mapping is clearly explained. The block-triangular Jacobian
    factorization for $\lambda$ is a nice result.
  \item The injectivity section is well-written but quite detailed for the main
    text. The proposition (lines 781--790) is the key result; the surrounding
    discussion of the compensation mechanism could be condensed by 30--40\%
    with more material pushed to Appendix~C (which already contains the full
    proofs).
  \item The proposition is presented in a \verb|\begin{quote}| with
    \verb|\textbf{Proposition}|. This is fine stylistically, but some journals
    prefer a formal theorem environment. Check JoE style guidelines.
\end{enumerate}

\subsubsection{Prior specification (\S4.4, lines 800--915)}

\begin{enumerate}[label=(\alph*)]
  \item The pythontex integration for displaying hyperprior values is excellent
    for reproducibility.
  \item The direct prior (equation~\ref{eq:d_prior}) specifies priors on
    $r$, $\ln\phi$, $\ln\eta$, $\ln\lambda$. The log transforms are mentioned
    but the rationale (positivity constraints on $\phi$, $\eta$, $\lambda$) is
    not stated here. A brief parenthetical would help.
  \item The NIG mapping (equations 10--13) is the core application of the
    variance decomposition. The derivation of $\alpha_0 = 2 + 1/\kappa^2$ from
    the IG variance formula is given inline. This is correct and clear.
  \item The diffuse hyperprior from Gao et al.\ (2022) uses
    IG(0.001, 0.001). Given that the paper has discussed the Gelman (2006)
    pathology of IG($\varepsilon$, $\varepsilon$) priors, this comparison is
    well-motivated. However, the paper does not explicitly note the irony: the
    uninformative comparison prior is precisely the type of prior the first camp
    identifies as problematic. Making this connection explicit would strengthen
    the narrative.
\end{enumerate}

\subsubsection{Results (\S4.5--4.7, lines 925--1078)}

\begin{enumerate}[label=(\alph*)]
  \item The hyperparameter posterior discussion (lines 953--979) is substantive
    and nuanced. The observation about $\eta$ --- where hyperparameter
    posteriors differ but the unconditional distributions are similar, explained
    through the variance decomposition --- is a genuine insight that
    demonstrates the interpretive value of the $\rho$ framework.
  \item The CE posterior predictive section (lines 981--1037) effectively
    demonstrates the practical stakes: the diffuse prior puts
    substantial mass at extreme risk aversion and risk seeking.
  \item The sensitivity table (\S4.7, lines 1041--1078) demonstrates the
    key invariance property: the prior predictive is stable across the
    $(\rho, \kappa)$ grid. The interpretation (lines 1070--1077) that the
    differences emerge in the \emph{posterior} is important and well-stated.
  \item \textbf{Missing}: A sensitivity analysis showing how the \emph{posterior}
    changes with $(\rho, \kappa)$ would be valuable. The prior predictive table
    demonstrates that the methodology works as advertised, but the reader's
    natural question is: ``so does $(\rho, \kappa)$ actually matter?'' A table
    or figure showing posterior summaries across the $(\rho, \kappa)$ grid
    would answer this. This could be computationally expensive (requiring
    MCMC for each grid point) or could be approximated analytically.
  \item Figure captions are terse (e.g., ``Hyperparameter posteriors''). JoE
    typically expects self-contained captions that summarize the key message.
    The \verb|\mycaption| command is used; check that its expansion provides
    adequate detail.
\end{enumerate}


\subsection{Sections 5 and 6: CPT and Discounting (lines 1081--1137)}

These are stubs. The outlines are reasonable: multiple DGPs, sample size grids,
100+ replications, three estimation approaches per replication. The scope is
ambitious.

\begin{enumerate}[label=(\alph*)]
  \item The stub text is currently in \verb|\textit{}|. This is fine for
    internal drafts but should be removed or replaced before any external
    review.
  \item Consider whether \emph{both} simulation studies are necessary for the
    first submission. A single well-executed simulation (CPT, given its higher
    parameter count and stronger motivation) might be sufficient, with
    discounting reserved for a revision or supplement. This would substantially
    reduce the time to submission.
  \item If both simulations are retained, the discounting section should clarify
    why it adds value beyond CPT. The current stub mentions that indifference
    amounts are ``directly elicited,'' which is a good distinguishing feature.
\end{enumerate}


\subsection{Section 7: Discussion (lines 1139--1213)}

\begin{enumerate}[label=(\alph*)]
  \item[{[x]}] The practical guidance (\S7.1) workflow labels have been changed
    from \verb|\textbf{}| to \verb|\emph{}|. The \verb|enumerate| environment
    is retained for the step-by-step procedure.
  \item The Bayes factor discussion (\S7.3, lines 1198--1213) is interesting
    but feels incomplete. The argument that BFs create a ``perverse incentive to
    retroactively tune priors'' is strong but needs a citation or a more careful
    statement. The point is valid for \emph{marginal likelihoods} used as model
    comparison tools; it is not unique to the framework proposed here.
    Consider citing \textcite{Gelman2013} or the ``garden of forking paths''
    literature. Alternatively, this subsection could be shortened to
    acknowledge the issue without fully resolving it.
  \item The limitations section (\S7.2) is honest and well-calibrated. The three
    limitations (moment-based, injectivity requirement, no model
    misspecification) are the right ones to highlight.
\end{enumerate}


\subsection{Section 8: Conclusion (lines 1215--1244)}

The conclusion is concise and effective. The ``reduces further'' phrasing
(lines 1233--1236) has been reworded to clarify that sharing $(\rho, \kappa)$
means only 2 of the $2p+2$ choices involve population-level reasoning, without
falsely implying a further numerical reduction.


\subsection{Appendix A: Full Covariance Derivation (lines 1249--1378)}

Complete and well-structured. The derivation is clean and the verification
confirms the algebra. A footnote has been added distinguishing the NIW scaling
parameter $\kappa_0$ from the CV of heterogeneity $\kappa$.

\subsection{Appendix B: Skew-Normal Extension (lines 1381--1456)}

Complete. Demonstrates the generalization to 3-moment matching. No issues.

Note: the \verb|\begin{enumerate}| on line 1401 for the three terms of the law
of total cumulance technically violates the no-bullet-point rule, but numbered
lists of mathematical terms seem acceptable.

\subsection{Appendix C: Global Injectivity (lines 1459--1557)}

Complete and detailed. The numerical verification, counterexamples, and
sufficient conditions are all present. Bold text has been removed (line 1534)
and sufficient condition labels changed to \verb|\emph{}|. One remaining
suggestion: the appendix is thorough enough that the main text injectivity
discussion (\S4.3) could be further condensed with a forward reference.


\subsection{Appendix D: Jacobian Correction (lines 1561--1619)}

Clear and correct. The three-part structure (why only forward map, geometric
intuition, non-injectivity) is pedagogically effective. No issues.

\subsection{Appendix E: MCMC Implementation (lines 1624--1707)}

Complete. The coupled transform sampler and NCP descriptions are clear. One
note: the ``115 blocks'' count is hardcoded (line 1639, 1645) rather than
computed from pythontex. If the number of subjects or blocking strategy changes,
these would need manual updates. Consider using \verb|\py{}| here too.

\subsection{Appendix F: Simulation Study Details (lines 1711--1721)}

Stub. Will be populated when Sections 5--6 are written.


\section{Cross-Cutting Issues}

\subsection{Notation}

\begin{enumerate}[label=(\alph*)]
  \item[{[x]}] $\kappa_0$ vs $\kappa$ confusion resolved via footnote in
    Appendix~A.
  \item The paper uses both $\bv{}$ (\verb|\bv|, $\boldsymbol{\vartheta}$) and
    $\bth$ (\verb|\bth|, $\boldsymbol{\theta}$). The distinction (generated
    quantities vs.\ structural parameters) is clear once established, but
    $\vartheta$ and $\theta$ are visually similar in some fonts. Monitor how
    this renders in the JoE template.
  \item The Prelec PWF uses $c_i$ for cumulative probabilities (line 635) and
    the law of total cumulance footnote uses $c_n$ for cumulants (line 516).
    These are in different sections and unlikely to confuse, but worth noting.
\end{enumerate}

\subsection{First-person consistency --- RESOLVED}

Fixed: ``we will only focus'' (line 421) and ``we are primarily concerned''
(line 450) changed to first person singular. Remaining ``we'' instances
(lines 496, 1427, 1429) are in mathematical exposition where the convention
is acceptable.

\subsection{Bold text --- RESOLVED}

All \verb|\textbf{}| in the paper body has been replaced with \verb|\emph{}|
or removed. Remaining instances are only in front matter
(\verb|\textbf{Keywords:}| and \verb|\textbf{JEL Classification:}|), which is
standard.

\subsection{Figure captions}

All figure captions use \verb|\mycaption{}| and appear terse from the LaTeX
source. Check that this command expands to include sufficient detail for
standalone interpretation.


\section{Prioritized Recommendations}

Based on the assessment above, I recommend the following order of operations,
organized by priority.

\subsection{Priority 1: Fixes that can be made immediately (no computation)
--- COMPLETE}

All items in this priority have been implemented:

\begin{enumerate}
  \item[{[x]}] Fixed ``an quantity-space'' $\to$ ``a quantity-space'' (line 167)
  \item[{[x]}] Fixed ``we'' $\to$ ``I'' (lines 421, 450)
  \item[{[x]}] Replaced \verb|\textbf{}| with \verb|\emph{}| in Discussion
    workflow labels (5 items), Proposition (\S4.3), and Appendix~C sufficient
    conditions
  \item[{[x]}] Removed bold from Appendix~C numerical verification (line 1536)
  \item[{[x]}] Added footnote distinguishing $\kappa_0$ (NIW) from $\kappa$
    (CV) in Appendix~A
  \item[{[x]}] Rewrote conclusion's ``reduces further'' phrasing to clarify
    that sharing $(\rho, \kappa)$ means only 2 of $2p+2$ choices involve
    population structure
  \item[{[x]}] Added sentence explaining log transforms ($\phi$, $\eta$,
    $\lambda$ are strictly positive) before direct prior interpretation
  \item[{[x]}] Made IG($\varepsilon$, $\varepsilon$) connection explicit in the
    diffuse prior comparison, linking back to Gelman (2006)
\end{enumerate}

\subsection{Priority 2: Structural improvements (editing, no computation)}

\begin{enumerate}
  \item Condense the belief-elicitation literature in \S2.2 (lines 303--338)
    to 2--3 sentences or a footnote. This tightens the paper by approximately
    1 page.
  \item Condense the main-text injectivity discussion (\S4.3) by moving the
    compensation mechanism details to Appendix~C. Keep the proposition and the
    practical recommendation (``use at least two distinct prize vectors'').
  \item Consider splitting Section~3 into two sections (one per contribution)
    to match the ``two contributions'' framing. This is optional but would
    improve structural clarity.
  \item Strengthen the Bayes factor discussion or shorten it. The current
    argument about ``perverse incentives'' needs either a citation or more
    careful hedging.
\end{enumerate}

\subsection{Priority 3: Decide simulation study scope}

This is the most consequential decision for time-to-submission.

\begin{enumerate}
  \item \textbf{Option A: Submit with one simulation (CPT only).} Drop the
    discounting section or reduce it to a brief discussion of how the framework
    applies. The CPT study is the stronger motivating example (highest
    parameter count, strongest $4p$ problem). The discounting study could be
    added in revision if requested.
  \item \textbf{Option B: Submit with both simulations.} Retain the current
    structure. This requires implementing two new model posteriors, simulation
    DGP code, and running hundreds of MCMC replications for each.
  \item \textbf{Option C: Submit with real-data application only.} Drop both
    simulation stubs and reframe as a methods paper with a single worked
    application. This would require strengthening the sensitivity analysis
    (Priority~4) to compensate for the lack of controlled evidence.
\end{enumerate}

My recommendation is \textbf{Option A}. The CPT simulation provides controlled
evidence that the methodology works, the RDU application provides real-data
credibility, and the paper would be a manageable length (~30 pages with
appendices). The discounting application can be added if referees request it.

\subsection{Priority 4: Posterior sensitivity analysis (computation required)}

The current sensitivity table (Table~1) shows that the \emph{prior predictive}
is invariant to $(\rho, \kappa)$. This is by construction and therefore
unsurprising. The more informative question is: how does the \emph{posterior}
vary with $(\rho, \kappa)$?

Two approaches:

\begin{enumerate}
  \item \textbf{Full MCMC grid:} Re-estimate the model at each $(\rho,
    \kappa)$ grid point. Computationally expensive but definitive.
  \item \textbf{Analytical approximation:} Use the fact that the posterior mean
    in a Normal--Normal model is a weighted average of the prior mean and the
    data mean, with weights determined by the prior precision and the data
    precision. This would give approximate posterior sensitivity without
    re-running MCMC.
\end{enumerate}

Either approach would substantially strengthen the paper. The current table
demonstrates correctness; a posterior table would demonstrate relevance.

\subsection{Priority 5: Phase 3 completion (CE-space priors)}

The CE-space prior specification (\S4.2) is described in the paper but
Phase~3 implementation is incomplete. Completing this would allow a
comparison of parameter-space vs.\ CE-space priors in the results section,
demonstrating Contribution~1 empirically (currently only Contribution~2 is
demonstrated).

This is important but can be done in parallel with the text improvements.

\subsection{Priority 6: Phase 4 (choice probability prior for $\lambda$)}

Currently $\lambda$ is handled with a direct parameter-space prior, not via
the quantity-space framework. Implementing the choice-probability mapping
would make the application fully consistent with the methodology. This is
a moderate-effort task that could be deferred to revision.


\section{Recommended Order of Operations}

\begin{enumerate}
  \item[{[x]}] Immediate text fixes (Priority 1) --- complete
  \item Structural editing (Priority 2): half a day
  \item Decide simulation scope (Priority 3): decision point
  \item Posterior sensitivity analysis (Priority 4): depends on approach
  \item CE-space implementation (Priority 5): multi-day
  \item Choice probability prior (Priority 6): multi-day, can defer
  \item Simulation study (Priority 3 execution): weeks
  \item Final polish (Phase 6): after all above
\end{enumerate}

The key decision point is Priority 3. Everything else follows from that
choice. If Option~A is selected, the path to a submittable draft requires
Priority~2 (text), Priority~4 (posterior sensitivity), one simulation
study (CPT), and final polish.


\section{Summary of Issues Found}

\begin{center}
\begin{tabular}{lllll}
\toprule
Location & Type & Severity & Status & Description \\
\midrule
Line 167 & Typo & Minor & Fixed & ``an quantity-space'' \\
Lines 421, 450 & Style & Minor & Fixed & ``we'' $\to$ ``I'' \\
\S2.2 (303--338) & Structure & Moderate & Open & Belief-elicitation detour \\
\S3 title & Structure & Minor & Open & Title doesn't reflect both contributions \\
\S4.3 & Length & Minor & Open & Injectivity could be condensed \\
\S4.4 line 858 & Clarity & Minor & Fixed & Log transforms unexplained \\
\S4.4 line 906 & Narrative & Minor & Fixed & IG($\varepsilon$,$\varepsilon$) irony \\
\S4.7 & Missing & Moderate & Open & No posterior sensitivity analysis \\
\S7.1 & Formatting & Minor & Fixed & Bold in workflow labels \\
\S7.3 & Argument & Minor & Open & BF ``perverse incentive'' needs citation \\
\S8 lines 1233--36 & Logic & Minor & Fixed & ``reduces further'' phrasing \\
App.\ A & Notation & Minor & Fixed & $\kappa_0$ vs $\kappa$ footnote \\
App.\ C & Formatting & Minor & Fixed & Bold text \\
App.\ E & Maintenance & Minor & Open & Hardcoded ``115 blocks'' \\
\S5--6 & Completeness & Major & Open & Stub sections \\
\bottomrule
\end{tabular}
\end{center}


\end{document}
